% ==============================================================================
% CHƯƠNG 4: KẾT QUẢ VÀ THẢO LUẬN
% ==============================================================================
\chapter{Kết quả và thảo luận}
\label{chap:ket-qua}

\section{Kết quả định lượng}
\label{sec:dinh-luong}

Các con số quan trọng nhất có thể đặt trong ô chỉ số để người đọc nắm ngay:

\begin{hangchiso}[3]
  \ChiSo{16\%}{tăng số tin tuyển dụng AI khi mức tiếp xúc tăng 1 độ lệch chuẩn}
  \ChiSo[do]{$-$12\%}{giảm tin tuyển dụng ở vị trí phi AI}
  \ChiSo[xanhla]{8,8 lần}{số tin yêu cầu kỹ năng AI, giai đoạn 2010--2018}
\end{hangchiso}

\Bang{tab:tong-hop} tổng hợp các chỉ số chính.

\begin{bangbieu}{Tổng hợp các chỉ số định lượng chính}{tab:tong-hop}{C{1cm} Y C{2.6cm} C{2.6cm}}
  \hangtieude \thead{STT} & \thead{Chỉ số} & \thead{Hệ số} & \thead{Sai số chuẩn} \\
  \midrule
  1 & Tin tuyển dụng AI            & \otot 15,96   & 1,73 \\
  2 & Tin tuyển dụng phi AI        & \oxau $-13{,}80$ & 4,22 \\
  3 & Việc làm cấp ngành           & $-0{,}05$     & 0,08 \\
  4 & Tiền lương theo nghề         & \onoibat $-0{,}17$ & 0,06 \\
\end{bangbieu}
\nguonbang{Tổng hợp từ \cite{acemoglu2022}.}

\section{So sánh và thảo luận}
\label{sec:thao-luan}

\SoSanh{Điểm mạnh}{
  \begin{dsdung}
    \item Dữ liệu bao phủ gần toàn bộ thị trường.
    \item Ba thước đo độc lập cho kết quả nhất quán.
  \end{dsdung}}
{Hạn chế}{
  \begin{dssai}
    \item Dữ liệu dừng trước giai đoạn AI tạo sinh.
    \item Chỉ phản ánh thị trường Hoa Kỳ.
  \end{dssai}}

\Bang{tab:so-sanh} là ví dụ bảng sọc, dễ đọc khi có nhiều hàng.

\begin{bangsoc}{So sánh các nghiên cứu được phân tích}{tab:so-sanh}{L{3.4cm} Y C{2.2cm}}
  \hangtieude \thead{Nghiên cứu} & \thead{Kết quả chính} & \thead{Nguồn} \\
  \midrule
  Acemoglu và cs. (2022)  & Tuyển dụng phi AI giảm tại cơ sở tiếp xúc cao & \cite{acemoglu2022} \\
  Brynjolfsson và cs. (2025) & Năng suất tăng 15\%, mạnh nhất ở nhóm ít kinh nghiệm & \cite{brynjolfsson2025} \\
  Acemoglu \& Restrepo (2020) & Mỗi robot/1.000 lao động giảm 0,2 điểm \% việc làm & \cite{acemoglu2020} \\
  Nghiên cứu khác & \canbosung{điền kết quả} & -- \\
\end{bangsoc}

\begin{canhbao}[Diễn giải thận trọng]
  Hệ số không có ý nghĩa thống kê \quantrong{không} đồng nghĩa với ``không có tác động''.
\end{canhbao}
